\section{A laboratory protocol}\label{sec:protocol}

The certificates of Section~\ref{to:sec:measure} are statements about any device under the contract. A laboratory test needs more: a concrete apparatus, a declared physical interface, a way to certify the budgets the contract assumes, and controls that separate the order price from everything else. This section designs such a test for the classical witness, the task of Figure~\ref{fig:three}, in a colloidal feedback trap of the kind used to test Landauer's principle~\cite{Jun14}; an earlier test used a modulated optical double well~\cite{Berut12}, and a preferred design with a preprogrammed physical memory is given below. Nothing quantum is needed, so the comparison with the quantum certificate of Proposition~\ref{to:prop:certificate} is a matter of platform, not of principle. That access to a copy changes the cost of erasure is established in theory~\cite{Bennett73,Bennett03,delRio11} and was measured in an electronic memory~\cite{Orlov12} (Section~\ref{sec:related}). What this paper contributes is the general schedule law and its consequences; what the proposed test adds is a controlled release schedule with control arms, a graded helper control and an explicit ledger of records and blank resources. The bound the test targets holds for every implementation, and its numbers are below. Two things it assumes are not yet established for a real apparatus: the record-export interface and the bounds on the controller in the ledger. Until they are, the protocol is a design proposal rather than a closed whole-device test.

\paragraph{The apparatus.} A bead in a double-well potential stores one bit; a second double well, or a second bead, stores the second. The potentials are shaped by feedback electronics that read the bead position with a camera and update the trap. The observer is a random number generator and a logging computer that supplies $a$, later $b$, and records every value the device returns. The feedback-trap test of Landauer's principle reads the work done on the bead by the trap from its trajectory~\cite{Jun14}, with the work and heat of stochastic energetics~\cite{Sekimoto98}, and work is not heat. For each episode we infer the heat into the bead's bath, at room temperature, from $Q=W_{\rm on}-\Delta U$, using trajectory work and a measured or bounded endpoint energy change. Mean work equals mean heat only after establishing equal mean initial and final bead energies for the complete cycle. We include every reset stroke and account separately for energy exchanges involving the other components of the declared interface.

The endpoint rule is part of the preregistered protocol. Each bead's endpoints are taken at lawful cuts: a carrier's when it arrives and when it is released, and an internal memory's before and after its accounted reset. A released carrier is not touched again except as a declared loan, and no endpoint energy is recorded or exported during service. At every cut the bead is in equilibrium within its logical well, not across the barrier, which would make the bit fair: given the well, its position has the calibrated canonical distribution of one fixed endpoint potential, symmetric under the exchange of its two wells, and transitions between the wells during equilibration are rare and counted as logical errors. The reflection symmetry then makes the mean energy the same in both wells, as in the cyclic erasure protocols of~\cite{Gavrilov17}; a symmetric potential alone would not, since the distributions within the two wells could differ. Departures from local equilibrium are bounded for each arm separately, from calibration runs separate from the data, and that bound is the residual mean $\Delta U$, whose uncertainty is part of the quoted precision of the mean heat. The central READY well of the preferred design below has an endpoint Hamiltonian of its own and needs its own energy correction. Positions suffice for the bead, which is overdamped: its momentum relaxes far faster than any stroke, so its kinetic energy is at equilibrium at every cut. The underdamped mechanical option below needs the rule in phase space.

\paragraph{Three arms.} Each episode runs one of three schedules, chosen at random so that drifts affect all arms alike.
\begin{enumerate}[label=(\roman*),leftmargin=*]
\item \emph{Held.} The device receives $a$, keeps it in the first well, receives $b$ in the second well, computes $a\oplus b$ into the second well by a two-well gate (a conditional barrier tilt), and returns both wells. Nothing is erased, so an ideal implementation needs no heat in the slow limit.
\item \emph{Lent.} The device receives $a$, copies it into a third well, hands the first well to the observer, who reads it and hands it back unmodified, receives $b$, computes $a\oplus b$, uncomputes the copy against the returned $a$ by the inverse gate, and returns everything. An ideal implementation again needs no heat in the slow limit. This arm is the control for early delivery without permanent release.
\item \emph{Early.} As in (ii), but the observer keeps the first well. The device computes $a\oplus b$ from its copy, returns the second well, and must then reset the copy, which costs Landauer's $\kB T\ln2$ at any speed.
\end{enumerate}
Arms (i) and (iii) realize the two bars of Figure~\ref{fig:three}; arm (ii) tests that the price is in the release and not in the delivery.

\paragraph{The quantity.} For each arm measure, per episode, the heat $Q$ into the bath and the disagreement $\delta$ between the first returned value and $(\text{second returned value})\oplus b$, which the observer computes from its own records. Proposition~\ref{to:prop:certificate}(d), in its classical form (Appendix~\ref{to:app:rank}), bounds the early arm from below: on fresh independent uniform inputs every episode must dispose of at least $1-h_2(\delta)$ bits of entropy. The probe source is part of the declared interface. It is independent of the device, and its freshness is conditional: given the arm choice and everything that happened before the episode, $a$ is a uniform bit independent of the device's state at the start of the episode, and $b$ is a uniform bit independent of $a$ and of the device's state when $b$ arrives; the helper noise of the graded control below is fresh in the same sense. Fair marginals are not enough: if every episode received the same fair bit, an exact device that kept a copy of each $a$ would hold one bit of record in all rather than one per episode, with every input fair and $\delta=0$. A source with a known bias or correlation needs its own entropy correction. A cyclic device restores its own memory, so this entropy goes into the bath or into the non-bath resources that are not restored: the controller, the battery and every other sink inside the declared interface. The heat into the bath, plus $\kB T\ln2$ times the entropy those resources absorb, is therefore at least $(1-h_2(\delta))\,\kB T\ln2$, while ideal held and lent implementations need no heat in the slow limit. With $\delta=5\%$ the bound is $0.71\,\kB T\ln2$; with $\delta=1\%$ it is $0.92$.

\paragraph{The interface.} The bound is a statement about a declared physical interface, which the protocol must fix before any data are taken. In the baseline three-arm protocol the only data carriers are two wells: the first arrives holding $a$, and the second arrives holding $b$ and is released holding $a\oplus b$. The declared schedule of each arm fixes every arrival and every release; in arm~(ii) the first well is, in addition, returned on a second visit and released again. The graded control below adds one declared visit of a helper well. Every other supplied blank, exported record or meter interaction must be included explicitly in the interface and its bound. Any other register that receives a record of a bead's position lies inside the boundary, as part of the controller: the camera and its frame buffer, the feedback registers and timers, the log of trap updates, and any computer that reads a frame, including the one that computes the heat. A blank log supplied to the device and a record exported from it are additional carrier visits, and placing them on the observer's side of the bookkeeping does not remove them. An exact reversible device shows why. Before releasing $a$, it copies $a$ into a blank memory $M$; when $b$ arrives it adds $M$ into the second well; then it swaps $M$ with a supplied blank log $L$ and exports $L$. On the well holding $a$, the second well, $M$ and $L$ this is a permutation, which sends $(a,b,0,0)$ to $(a,a\oplus b,0,a)$, so at degenerate Hamiltonian it runs the early arm with no heat, a blank memory and $\delta=0$. The record has left through an extra output, using an extra blank input; the bound is not violated, because its interface has changed. Two more loopholes fix what the audit must cover. First, a record need not be a record of $a$ to carry $a$ away. If the device instead copies $b$ into $L$ when $b$ arrives, adds $M$ into the second well, and then clears $M$ by adding the second well and $L$ into it, the same registers go from $(a,b,0,0)$ to $(a,a\oplus b,0,b)$, and on the carriers alone this is the reversible map $(a,b,0)\mapsto(a,a\oplus b,b)$. The exported record of $b$ carries no information about $a$ by itself, but reveals $a$ jointly with the output $a\oplus b$, so joint records must be audited, not each record on its own. Second, a supplied blank need not leave with any record to pay for the reset. Letting an independent blank bit relax to a random state can return $\kB T\ln2$ of work, since entropy can be exported by randomizing degrees of freedom outside the data~\cite{Bennett03}; that work pays for the reset of the copy, and the blank leaves random and independent of $a$. Every consumed blank resource therefore counts, not only records of $a$. For the two-carrier bound used here, no record leaves the boundary before the accounted reset of the episode, and every blank resource the episode consumes enters the ledger. At that reset every register inside the boundary is restored to a fixed state, using as side information only what the device still holds (in arm~(ii), the returned well while it is held again), never a released well or the observer's log, and the dissipation of each restoration is counted in the ledger below. The heat itself is computed from the bead's trajectory, which is such a record, so its readout is a record export too. A closed test must therefore include the logging hardware and its restoration in a declared device boundary and prove the corresponding full-interface bound; we have not done so.

\paragraph{The ledger: what the controller may not hide.} The lower bound charges the sum of the entropy changes of every sink inside the interface. In a feedback trap the dangerous sink is the controller itself: the camera frames, the feedback registers and the log of the trap updates are records, and a record of $a$ that survives the episode is a stranded copy that was never erased, which makes the measured heat an underestimate by exactly the amount the bound is testing. The protocol therefore keeps a ledger with four entries. They combine measured quantities with explicit, independently justified bounds on the degrees of freedom that are not observed; calibrating the trap parameters alone does not certify the entropy of the whole controller.
\begin{enumerate}[label=(\alph*),leftmargin=*]
\item \emph{Camera and feedback records.} The feedback loop is run from a fixed program whose state after the episode is restored to a fixed value; every frame used to drive the trap is overwritten by the next, and the buffer is cleared at the end of each episode. The dissipation of every overwrite and of the clearing is counted: in $Q$ if it goes into the bead's bath, and otherwise as the entropy change of the sink that receives it, at that sink's own temperature. No frame leaves the boundary before that reset.
\item \emph{The controller.} The potential parameters at the start and at the end of each episode are compared. Any difference is a change of the controller's state; its entropy and energy are estimated from the parameter distribution over episodes and entered in the ledger as $\Delta S_C$ and $\Delta E_C$. A reproducible control voltage does not show that the whole controller returned: the electronics, hidden buffers, timers, software state and the correlations among them need bounds of their own. Present feedback electronics also dissipate far more than $\kB T$ per operation into their surroundings. The bound counts those surroundings among the other sinks, so it says nothing about the heat into the bead's bath alone unless a sharper accounting isolates the part of their entropy that comes from the device's records; measuring their heat does not do that. The feedback-trap test of Landauer's principle accounts for the work done on the bead, not for the camera, the computer or the illumination~\cite{Jun14}; these bounds are what this design adds to it, and they are not yet established.
\item \emph{The battery.} The source that drives the trap is the battery. Its entropy change $\Delta S_W$ is a separate entry and needs a bound of the same kind.
\item \emph{The bath.} Use the measured initial bath temperature for each analyzed block, and bound any initial departure from equilibrium and any external energy exchange explicitly. Lemma~\ref{lem:identity} assumes a bath that starts in its Gibbs state. For a bath that starts in $\rho_{R,i}$, with $\gamma_\beta$ the Gibbs state at the reference temperature and entropies in nats, $\beta Q=\Delta S_R+D(\rho_{R,f}\|\gamma_\beta)-D(\rho_{R,i}\|\gamma_\beta)$, and the initial relative entropy subtracts. A drift of the starting temperature, or an unmonitored exchange of heat, therefore cannot be entered as a nonnegative final correction; the final relative entropy may be dropped only once the initially Gibbs assumption at that temperature is justified.
\end{enumerate}
The reported quantity is $Q+\kB T\ln2\,(\Delta S_C+\Delta S_W+\Delta S_{\rm other})$, per episode and per arm, with every term referring to the same declared implementation and $\Delta S_{\rm other}$ covering every sink inside the interface not named above, every consumed blank resource included. The early arm must exceed $(1-h_2(\delta))\kB T\ln2$ at every stroke duration, whatever overhead the apparatus adds. The ideal held and lent implementations need no heat in the slow limit, but accurate accounting does not make an inefficient apparatus approach that optimum. A held episode whose camera keeps a frame showing $a$ until both outputs have left must reset that frame without the released carriers, which on fresh uniform inputs disposes of one bit however slowly the gates run. We have not shown that the complete feedback apparatus, with its readout and record restoration, approaches the optimum. A test of the held and lent intercepts therefore also needs a construction that cleans these records while the device still holds their side information, or an independently justified bound on the residual overhead. The difference between the early arm and the held arm at the same $\delta$ is the measured excess of these two implementations. It equals the order price only if both implementations attain their respective optima under the same budgets; the bound on the early arm needs no such optimality. The lent arm, run near its optimum, checks that early delivery by itself costs nothing.

\paragraph{Numbers.} At room temperature ($300$\,K), $\kB T\ln2\approx2.9\times10^{-21}$\,J, the scale resolved in~\cite{Berut12,Jun14}. Confidence in the error and confidence in the heat are separate questions. The disagreement indicators are bounded, so the Hoeffding--Azuma inequality~\cite{Hoeffding63,Azuma67}, which holds when the device carries memory across episodes, applies to them: with $5\times10^4$ episodes per arm the observed disagreement is within $10^{-2}$ of the average conditional disagreement probability, except with probability below $10^{-3}$. That argument does not carry over to the heat, whose values are unbounded and may be correlated across episodes; a shared offset, such as a drift of the calibration, survives any amount of averaging. Heat confidence is therefore assessed on its own, from independence or mixing assumptions that the data must justify, or from an explicit covariance or bounded-increment analysis, with the calibration drift included. Under independence the uncertainty in the mean heat is $\sigma_Q/\sqrt N$. As a feasibility estimate, suppose that every stroke of an episode has a standard deviation of at most $2.2\,\kB T$, comparable to the fastest erasure cycles of Jun et al., whose statistical error of about $0.10\,\kB T$ after $60$ minutes of $7.5$-second cycles means a spread of about $2.2\,\kB T$ per cycle~\cite{Jun14}, and that strokes and episodes are independent. Ten strokes then give an uncertainty below $0.032\,\kB T$ at $N=5\times10^4$. Ten perfectly correlated strokes have ten times, not $\sqrt{10}$ times, the spread of one, which would put the uncertainty near $0.1\,\kB T$. The variance bound is an assumption, to be established for the copying, conditional-gate and restoration strokes of this apparatus; the erasure experiment of~\cite{Jun14} does not establish it for them, and independence alone does not supply it, since ten independent strokes of spread $100\,\kB T$ would give $1.4\,\kB T$. The uncertainty of the endpoint correction adds to these figures. Their shortest erasure cycle, $7.5$\,s, is far from reversible: it costs about $3.5\,\kB T$ of work, five times $\kB T\ln2$, and their Landauer value, $0.71\pm0.03\,\kB T$, is the intercept of a sweep over cycle durations extrapolated along $1/\tau$~\cite{Jun14}. The protocol therefore runs every arm over such a duration sweep and compares the intercepts. The bound on the early arm holds at every duration, so a fast erasure that dissipates more is still a valid early arm; only the held and lent intercepts, whose ideal value is zero, need the extrapolation. At a $7.5$-second stroke time, $5\times10^4$ episodes require $4.34$ days times the number of strokes in an episode, per arm and per duration of the sweep: about $43$ days for ten strokes, before overhead, and longer durations take proportionally longer. Faster operation requires a separate precision and dissipation estimate.

\paragraph{Preregistration.} Before any data are taken, the protocol fixes the probe source and the randomization of the arms, the duration grid, the helper-noise grid of the graded control below, the number of episodes at every point, and the total failure probability with its split among the error, heat, endpoint and calibration estimates. At each point the claimed bound is its least value over a confidence interval for the disagreement. Since $1-h_2(\delta)$ and $h_2(\eta)-h_2(\delta)$ decrease in $\delta$ only up to $1/2$, it suffices to substitute $\delta^*=\min\{\delta_U,1/2\}$ for $\delta$ in both, with $\delta_U$ an upper confidence bound; $\delta_U$ itself would overstate the claim once it passes $1/2$. The noise parameter $\eta$ is set by the observer's random number generator and calibrated; if it is estimated instead, the claim uses the least value of $h_2(\eta)$ over its confidence interval. A claim about a whole sweep uses intervals that hold simultaneously at every point, for instance by a union bound over the points. The bounds concern expected values, and the Hoeffding--Azuma estimate concerns the run's own average conditional disagreement probability; the two agree when the episodes repeat independently, which the reset at the end of every episode is designed to ensure. Independent repetition, like the heat model with its bounds on drift, covariance and the variance of each stroke, is a preregistered assumption, not a consequence of the data.

\paragraph{The loan as a contract.} Arm (ii) lies outside the observer class of Theorem~\ref{to:thm:law}, whose observers may act on a loan before returning it. It is covered by a contract with one additional promise: the observer returns the lent carrier in the state it received it, jointly with everything else and not only in distribution. In this classical protocol the returned value must equal the delivered value in every episode, whatever else the episode records: a fresh fair bit returned in its place has the right distribution, but it is independent of the copy, which is then stranded as in arm~(iii). For a quantum carrier the promise is that the observer acts on the loan as the identity, including on its retained references. Under that promise the device may treat the loan as a held carrier for the purpose of cleanup, and the law applies to the schedule with the loan removed; the price of the task is then zero, which is what arm (ii) measures.

A deliberately flipped return is a cheap additional control. It tests how the declared implementation responds when the promise is broken, and what it shows depends on the circuit. Write the first well, the second well and the copy as $(A,B,M)$, starting from $(a,b,0)$. The device sets $M\leftarrow M\oplus A$ and lends $A$, and the observer returns it holding $a\oplus1$. A device that computes the answer from the returned well, $B\leftarrow B\oplus A$ and then $M\leftarrow M\oplus A$, ends in $(a\oplus1,a\oplus b\oplus1,1)$: every answer is wrong, and $\delta=1$. A device that computes it from the copy, $B\leftarrow B\oplus M$ before $M\leftarrow M\oplus A$, ends in $(a\oplus1,a\oplus b,1)$: every answer is right, and $\delta=0$ as defined above, from the first delivery of $A$, while a disagreement computed from its final release would be $1$. In both circuits the copy ends in the fixed state $1$, not in a record of $a$, and a fixed NOT restores it at no cost: the flipped well still carries the whole bit $a$, and the flip changes only the relation that the uncomputation expects. The control therefore records $\delta$, with the release of $A$ that it uses declared, and the state in which $M$ is left, read at the end of the episode as a declared meter interaction, and compares them with the declared circuit. It does not certify that no hidden record survives: an unused copy of $a$ kept elsewhere changes neither observable. That rests on the ledger and its independently justified bounds.

\paragraph{A graded control.} The lent and early arms are the two ends of a curve. In a variant of the early arm, after the first well has been released for good, the observer lends the device a helper well holding $R=a\oplus N_\eta$, where the noise bit $N_\eta$ equals $1$ with probability $\eta$ and is independent of everything else; the device returns the helper unmodified. The task outputs are unchanged and stay exact. Adding $R$ into the copy turns it into $N_\eta$, and a reset tuned to the distribution of $N_\eta$ then costs $h_2(\eta)\,\kB T\ln2$; a reset tuned to a fair bit would cost $\kB T\ln2$ at every $\eta$~\cite{KolchinskyWolpert17}. No device does better. Since $R$ returns unchanged, the helper can lower the device's entropy by at most the mutual information between the device and $R$, which is at most $I(a{:}R)=1-h_2(\eta)$ bits because, given $a$, the helper carries only fresh noise; with the $1-h_2(\delta)$ bits that the release forces, every device must dispose of at least $h_2(\eta)-h_2(\delta)$ bits per episode. With exact outputs the optimal cleanup heat is $h_2(\eta)\,\kB T\ln2$: $0$, $0.081$, $0.469$, $0.811$ and $1$ in units of $\kB T\ln2$ at $\eta=0$, $0.01$, $0.1$, $0.25$ and $0.5$. At $\eta=0$ this is the lent arm, and at $\eta=1/2$ the early arm. The control is a curve test of the access law: the price is the entropy of the stranded copy given what the device can still reach. Its cleanup is the erasure of a fraction of a bit, which Gavrilov, Ch\'etrite and Bechhoefer measured in a feedback trap: in the slow limit the work agreed with $\kB T\ln2$ times the Shannon entropy of the bit when the protocol was shaped to the initial distribution, while a protocol that first lowers the barrier cost $\kB T\ln2$ whatever the initial probability~\cite{Gavrilov17}. What the graded control adds is the source of the bias, a helper correlated with the released bit.

\paragraph{A preferred design.} The feedback trap makes the interface hard to close, because its controller records the bead's position at every update. A physical, preprogrammed memory avoids those records: each bit is a fixed double well driven by a program written before the episode, the drive never reads the data, and the early arm resets the copy by a symmetric merge into a central READY well. Under the reflection that exchanges the two wells, that reset is symmetric, and a meter whose complete coupling respects the symmetry has the same record distribution for $a=0$ and $a=1$; a tilted reset, as in the standard erasure protocols, gives no such guarantee, and strokes that are not symmetric need the property by design or an audit of their own. Two readouts can meet that condition: delayed calorimetry, with the thermometer disconnected during the task and read only after the accounted reset, for instance in superconducting flux logic with an engineered absorber; or, in a physical mechanical double well, an even-only readout of $x^2$, which determines the work of every stroke whose potential is even in $x$. That memory is underdamped, so $U$ includes kinetic energy: each endpoint must be locally canonical in phase space at the calibrated temperature, or the residual change of kinetic energy must be bounded independently, and the heat is $Q=W_{\rm on}-\Delta(E_{\rm pot}+E_{\rm kin})$, with the uncertainty of both terms in the heat interval. Experiments of Dago, Ciliberto and Bellon on such a memory, a micromechanical oscillator in a double well, show why this needs its own calibration: a fast erasure heats the oscillator, up to twice the bath temperature~\cite{DagoCilibertoBellon23}, and erasures repeated without time to thermalize warm it until, above a threshold speed, resets fail~\cite{DagoCilibertoBellon24}. They motivate this platform but do not close the interface required here. Joint records and consumed blanks still have to be audited as above, the meter's own entropy change enters the ledger, and this, too, is a design proposal.

\paragraph{What the test would show.} Once the interface and the ledger are closed, implementations that approach their respective optima would give three intercepts of the duration sweep: held near zero, lent near zero, and early above $(1-h_2(\delta))\kB T\ln2$ after the ledger, as is every early point of the sweep whatever the overhead; the graded control would add the curve $h_2(\eta)$ between the last two, never below its bound $h_2(\eta)-h_2(\delta)$. It would show, in one apparatus, a thermodynamic cost that is forced by a release schedule rather than by a map, and that survives arbitrarily slow operation. The early arm is a Landauer erasure of the stranded copy, so the experiment demonstrates the bound, which holds for every implementation, rather than testing it. A negative result, an early arm below the bound after the ledger, would mean that a premise of the bound failed: a copy of $a$ survived somewhere the ledger did not look, a record left through an undeclared export, a blank resource was consumed without being counted, the probes were not fresh, the bath did not start in its Gibbs state, or the sampling, the error calibration or the heat accounting was wrong. No control arm identifies which of these failed; a surviving copy in particular is excluded only by the ledger and its bounds.
